\scanpage{23}\mspage{43}
\[
H^4(X,\underline{\mathrm Z}_l\handwritten{[}2\handwritten{]})\handwritten{\simeq}\underline{\mathrm Z}_l\ .
\]
\deleted{D'après \unsure{Il est}} Cette forme \inserted{(en négligeant la torsion du $H^2$) est non dégénérée, et} indu\deleted{\unsure{1}}it sur $\mathrm{NS}(X)\otimes_{\underline{\mathrm Z}}\underline{\mathrm Z}_l$ la forme bilinéaire déduite de la \uline{forme intersection} sur $\mathrm{NS}(X)$, provenant de la \deleted{\unsure{p}}multiplicité globale d'intersections de diviseurs sur $X$ \margin{\handwritten{$\times$}}\handwritten{[15]} \deleted{\unsure{;}}. D'après \margin{\handwritten{$\times$}}MATSUSAKA \handwritten{[27]}, cette dernière forme est non dégénérée. T\deleted{\unsure{a}}ensorisant al\unsure{o}rs la suite exacte (8.7) par $\underline{\mathrm Q}_l$, on tro\deleted{\unsure{n}}uve que $T_l(\mathrm{Br}^{\deleted{\unsure{\prime}}}(X))\otimes_{\underline{\mathrm Z}_l}\underline{\mathrm Q}_l$ \uline{s'identifie canoniquement à l'orthogonal de} $\mathrm{NS}(X)\otimes_{\underline{\mathrm Z}_l}\underline{\mathrm Q}_l$ \uline{dans} $H^2(X,\underline{\mathrm Q}_l\handwritten{[}1\handwritten{]})$, \uline{et que le cup-produit définit sur cet esp\deleted{\unsure{e}}ace une forme quadratique non dégénérée}. L'interprétant comme une forme quadratique sur $T_l(\mathrm{Br}^{\deleted{\prime}}(X))$ à valeurs dans $\underline{\mathrm Q}_l$, (et tenant compte que le cup-produit induit même une autodualité de $H^2/\text{torsion}$\unsure{)|),|;)} on trouve aisément que cette forme est à valeurs dans $\underline{\mathrm Z}_l$ si et seulement si le discriminant de la forme intersection sur $\mathrm{NS}(X)/\text{torsion}$ est une unité $l$-adique - condition qui est vérifiée donc pour presque tous les nombre premiers $l$ \unsure{-}, et que dans ce cas on trouve une autodualité de $T_l(\mathrm{Br}^{\deleted{\prime}}(X))$.

\handwritten{8.3.} Ce résultat ne concerne essentiellement que la partie divisible $\mathrm{Br}'(X)^{\unsure{o|0}}$ du groupe de Brauer cohomologique. On peut également donner une relation de dualité concernant le \deleted{\unsure{grou}} groupe fini $\mathrm{Br}'(X)/\mathrm{Br}'(X)^o$, qui grâce à la suite exacte de Kummer (BR II 3.1) s'explicite comme
\[
\mathrm{Br}'(X)/\mathrm{Br}'(X)^o\handwritten{\simeq}H^2(X,\handwritten{\mu}_{l^\infty})/H^2(X,\handwritten{\mu}_{l^\infty})^o\ , \tag*{(8.8)}
\]
(compte tenu que $\mathrm{Pic}(X)\otimes\underline{\mathrm Q}_l/\underline{\mathrm Z}_l=\mathrm{NS}(X)\otimes\underline{\mathrm Q}_l/\underline{\mathrm Z}_l$ est divisible). Utilisant (8.3) pour $i=3$, on trouve donc
\[
\mathrm{Br}'(X)/\mathrm{Br}'(X)^o\handwritten{\simeq}\mathrm{Tors}\ H^3(X,\underline{\mathrm Z}_l\handwritten{[}1\handwritten{]})\ , \tag*{(8.9)}
\]
\handwritten{(}valable indépendamm\deleted{\unsure{a}}ent des hypothèses particulières \deleted{\unsure{faites}}\inserted{posées} au début de l'alinéa précédent pour $X$. Lorsque celles-ci sont vérifiées, alors les relations de dualité,bien connues dans le cas classique, et qui résultent aisément \inserted{\handwritten{par passage à la limite}} des résultats de dualité \margin{\handwritten{$\times$}}\handwritten{[}SGAA \handwritten{$\underline{\overline{\mathrm{XIX}}}$\quad]} pour des
% archive sheet number: 22
